\chapter{Introduction}

The introduction motivates the topic and provides an overview of the thesis.
It can be seen as a longer version of the abstract that follows the hourglass 
structure again, but with emphasis on the top part of the hourglass.
However, you can now use citations and explain everything in more detail.
It should be accessible to any computer scientist without prior knowledge of 
the specific area.

%-----------------------------------------------------------------------
% Background and Motivation
%-----------------------------------------------------------------------

The background and motivation part should cover the following points:
\begin{enumerate}
    \item Name the research area and describe informally what objects are
          studied (e.g., terms, integer programs, Java programs, automata) and
          what property is of interest (e.g., termination, complexity, safety).
    \item Explain why the property matters.
    \item Give a small, self-contained example that illustrates the problem informally.
    \item Explain why answering the question automatically is non-trivial and
          what makes the problem hard in general.
\end{enumerate}

%-----------------------------------------------------------------------
% The Specific Open Problem
%-----------------------------------------------------------------------

The open-problem part should contain:
\begin{enumerate}
    \item A summary of what is already known in the area (prior results,
          existing automated techniques, tool support).
    \item A precise identification of the gap: e.g., no automated procedure
          exists, existing techniques are unsound or incomplete for a relevant
          subclass, or no efficient algorithm is known.
    \item A clear statement of what this thesis aims to achieve.
\end{enumerate}

%-----------------------------------------------------------------------
% Related Work
%-----------------------------------------------------------------------

\subsection*{Related Work}

Discuss the most relevant prior work.
For each cited work, address the following:
\begin{enumerate}
    \item What does the work do? Describe the approach in one to three
          sentences and do not just list references.
    \item How does it relate to this thesis? Does it solve a special case,
          use a different technique, or target a different formalism?
    \item What limitation motivates your own contribution?
\end{enumerate}

%-----------------------------------------------------------------------
% Contributions
%-----------------------------------------------------------------------

\subsection*{Contribution}

List one item per concrete contribution. 
Each item should:
\begin{enumerate}
    \item Identify what is new (a technique, a criterion, a framework, an
          implementation, or an experimental evaluation).
    \item State the property proved, if applicable (soundness, completeness,
          both, or a complexity bound).
    \item Reference the theorem, or section in which the contribution appears.
\end{enumerate}

%-----------------------------------------------------------------------
% Thesis Structure
%-----------------------------------------------------------------------

\subsection*{Structure}

End the introduction by outlining the remaining chapters.
Give at least one sentence per chapter describing its content. 
Do not simply write ``Chapter~X presents the main results.'' 

Example: ``The remainder of this thesis is organized as follows.
\Cref{chap:preliminaries} introduces the formal background required throughout
this thesis, including \ldots
Then, \Cref{chap:main-chapter} presents \ldots
\Cref{chap:evaluation} describes the implementation and the experimental
evaluation.
Finally, \Cref{chap:conclusion} summarizes the results and discusses directions for
future work.''